%%%%%%%%%%%%%%%%%%%%%%%%%%%%%%%%%
%
%       EXERCÍCIO
%
%%%%%%%%%%%%%%%%%%%%%%%%%%%%%%%%%

\definevalorquestao

\begin{exercicioBanco}[\valorquestao]
Um reservatório contém 72 m\(^3\) de água e será drenado. O volume de água que sai é dado por \(V(t) = 24t - 2t^2\), onde \(t\), dado em horas, é o tempo decorrido após o início da drenagem. Se a dreanagem começou às 10 h, quando o reservatório estará vazio?

% Define as alternativas
\newcommand{\alternativas}{%
    \begin{center}
        \begin{tabularx}{\textwidth}{XXX}
            (a) \(22\) h. &
            (b) \(10\) h do dia seguinte. &
            (c) \(20\) h. \\[10pt]
            (d) \(19\) h. &
            (e) \(16\) h.
        \end{tabularx}
    \end{center}
}

% Define a resposta correta
\newcommand{\resposta}{E}

% Lógica condicional para exibição
\ifdefstring{\atividade}{lista}{%
    \alternativas
    \vspace{0.5em}
    
    \noindent\textbf{Resposta:} letra \textbf{\resposta}.
}{%
    \ifdefstring{\modo}{objetivo}{%
        \alternativas
    }{%
        % modo = discursiva → não mostra alternativas
    }
}
\end{exercicioBanco}

